% --- Archon physics formalization source begin ---
% archon:physics
% archon:covers PhyXMiniProblems/problem_phyx_mini_0648.lean
% archon:source-report reports/phyx_mini/problem_phyx_mini_0648.source.json

\chapter{Physics problem phyx\_mini\_0648}
\label{ch:PhyXMiniProblems_problem_phyx_mini_0648}

\paragraph{Problem source.}
\#\# Physical scenario

Figure shows the wave function of a neutron.

\#\# Question

What is the probability that the neutron is located between \$x = -1.0\$ \$mm\$ and \$x = 1.0\$ \$mm\$?

\#\# Answer choices

- A: A: 0.75

- B: B: 0.45

- C: C: 0.25

- D: D: 0.1

\#\# Figure caption (auxiliary; use the image as primary evidence)

The image is a graph with a piecewise function, plotted in red, on the Cartesian coordinate plane. 

- **Axes**: The horizontal axis is labeled \textbackslash{}(x \textbackslash{}, (\textbackslash{}text\{mm\})\textbackslash{}), and the vertical axis is labeled \textbackslash{}(\textbackslash{}psi(x)\textbackslash{}).
- **Horizontal Axis Range**: The x-axis ranges from -4 to 4.
- **Vertical Axis**: It is marked at \textbackslash{}(-c\textbackslash{}), 0, and \textbackslash{}(c\textbackslash{}).

**Function Description**:

- For \textbackslash{}(x < -4\textbackslash{}) and \textbackslash{}(x > 4\textbackslash{}), \textbackslash{}(\textbackslash{}psi(x)\textbackslash{}) is not defined or shown.
- From \textbackslash{}(x = -4\textbackslash{}) to \textbackslash{}(x = 4\textbackslash{}), the function alternates values:
  - \textbackslash{}(\textbackslash{}psi(x) = -c\textbackslash{}) from \textbackslash{}(x = -4\textbackslash{}) to \textbackslash{}(x = -2\textbackslash{}).
  - \textbackslash{}(\textbackslash{}psi(x) = 0\textbackslash{}) at \textbackslash{}(x = -2\textbackslash{}) and \textbackslash{}(x = 2\textbackslash{}).
  - \textbackslash{}(\textbackslash{}psi(x) = c\textbackslash{}) from \textbackslash{}(x = -2\textbackslash{}) to \textbackslash{}(x = 2\textbackslash{}).
  - \textbackslash{}(\textbackslash{}psi(x) = 0\textbackslash{}) at \textbackslash{}(x = 2\textbackslash{}) and \textbackslash{}(x = 4\textbackslash{}).
  - \textbackslash{}(\textbackslash{}psi(x) = -c\textbackslash{}) from \textbackslash{}(x = 2\textbackslash{}) to \textbackslash{}(x = 4\textbackslash{}).

This graph represents a step function with defined jumps at specific x-values, forming a symmetric pattern around the origin.

\#\# Dataset metadata

Quantum Phenomena; Numerical Reasoning, Physical Model Grounding Reasoning

\paragraph{Recorded answer/context.}
C: C: 0.25

\paragraph{Figure/image path.}
/root/proposal\_for\_physic/science-mango/phyx\_mini\_run/phyx\_data/test\_image/648.png

\paragraph{Formalization target.}
create a compiling Lean file with sorry bodies at `PhyXMiniProblems/problem\_phyx\_mini\_0648.lean`. The Lean declarations must preserve the physical quantities, dimensions or dimensional roles, figure labels, governing-law hypotheses, and final relation expressed by this problem.
Use Mathlib/Physlib names found through LeanExplore where available. If a physics API is missing, introduce faithful local abstractions rather than scalar placeholder aliases.

\begin{theorem}[Physics formalization target]
\label{thm:physics:phyx_mini_0648:target}
\lean{PhyXMiniProblems.ProblemPhyXMini0648.recorded_choice_is_correct}
\uses{def:physics:phyx-mini-0648:phyxminiproblems-problemphyxmini0648-neutronpositionexperiment, def:physics:phyx-mini-0648:phyxminiproblems-problemphyxmini0648-satisfiesonedimensionalpositionquantumlaws, def:physics:phyx-mini-0648:phyxminiproblems-problemphyxmini0648-matchesproblemstatementandfigure, def:physics:phyx-mini-0648:phyxminiproblems-problemphyxmini0648-requestedcentralinterval, def:physics:phyx-mini-0648:phyxminiproblems-problemphyxmini0648-matchesanswerchoice}
The assigned autoformalize agent should translate this physics problem into Lean declarations in the covered file, with theorem and lemma proof bodies written as `by sorry`.
\end{theorem}
\begin{proof}
This is an autoformalization task, not a proof task. Produce statements that can later be proved by the physics prover without weakening the physical model.
\end{proof}

% --- Archon declaration topology begin ---
\section{Declaration topology}
\begin{definition}[Quantum Particle Kind]
  \label{def:physics:phyx-mini-0648:phyxminiproblems-problemphyxmini0648-quantumparticlekind}
  \lean{PhyXMiniProblems.ProblemPhyXMini0648.QuantumParticleKind}
  The particle species named by the problem statement
\end{definition}
\begin{proof}
  This is the declared model definition of Quantum Particle Kind.
\end{proof}
\begin{definition}[Wave Function Curve Shape]
  \label{def:physics:phyx-mini-0648:phyxminiproblems-problemphyxmini0648-wavefunctioncurveshape}
  \lean{PhyXMiniProblems.ProblemPhyXMini0648.WaveFunctionCurveShape}
  Qualitative shape of the wavefunction in the supplied primary image
\end{definition}
\begin{proof}
  This is the declared model definition of Wave Function Curve Shape.
\end{proof}
\begin{definition}[Neutron Wave Function Figure]
  \label{def:physics:phyx-mini-0648:phyxminiproblems-problemphyxmini0648-neutronwavefunctionfigure}
  \lean{PhyXMiniProblems.ProblemPhyXMini0648.NeutronWaveFunctionFigure}
  \uses{def:physics:phyx-mini-0648:phyxminiproblems-problemphyxmini0648-wavefunctioncurveshape}
  Numerical and textual readouts from the plotted wavefunction. Coordinates are real readouts in millimetres. The amplitude scale is the real number printed as c; as a wavefunction amplitude it has the implicit unit mm⁻¹ᐟ², rather than being an untyped physical scalar
\end{definition}
\begin{proof}
  This is the declared model definition of Neutron Wave Function Figure.
\end{proof}
\begin{definition}[Neutron Position Experiment]
  \label{def:physics:phyx-mini-0648:phyxminiproblems-problemphyxmini0648-neutronpositionexperiment}
  \lean{PhyXMiniProblems.ProblemPhyXMini0648.NeutronPositionExperiment}
  \uses{def:physics:phyx-mini-0648:phyxminiproblems-problemphyxmini0648-quantumparticlekind, def:physics:phyx-mini-0648:phyxminiproblems-problemphyxmini0648-neutronwavefunctionfigure}
  A neutron position experiment with a coordinate representative of its quantum state and an independently specified probability observable. Normalization is carried by ProbabilityMeasure; Born's rule is imposed separately below
\end{definition}
\begin{proof}
  This is the declared model definition of Neutron Position Experiment.
\end{proof}
\begin{definition}[Satisfies One Dimensional Position Quantum Laws]
  \label{def:physics:phyx-mini-0648:phyxminiproblems-problemphyxmini0648-satisfiesonedimensionalpositionquantumlaws}
  \lean{PhyXMiniProblems.ProblemPhyXMini0648.SatisfiesOneDimensionalPositionQuantumLaws}
  \uses{def:physics:phyx-mini-0648:phyxminiproblems-problemphyxmini0648-neutronpositionexperiment}
  One-dimensional state admissibility and the position-space Born rule. Lebesgue volume is taken on the numerical millimetre coordinate. Thus Complex.normSq has reciprocal-millimetre density role, and withDensity turns it into the dimensionless position-probability measure
\end{definition}
\begin{proof}
  This is the declared model definition of Satisfies One Dimensional Position Quantum Laws.
\end{proof}
\begin{definition}[Matches Problem Statement And Figure]
  \label{def:physics:phyx-mini-0648:phyxminiproblems-problemphyxmini0648-matchesproblemstatementandfigure}
  \lean{PhyXMiniProblems.ProblemPhyXMini0648.MatchesProblemStatementAndFigure}
  \uses{def:physics:phyx-mini-0648:phyxminiproblems-problemphyxmini0648-neutronpositionexperiment}
  Literal labels, landmarks, and plateau values read from the primary image 648.png. This predicate contains no probability of the requested interval and no numerical answer-choice value
\end{definition}
\begin{proof}
  This is the declared model definition of Matches Problem Statement And Figure.
\end{proof}
\begin{definition}[Position Interval Readout]
  \label{def:physics:phyx-mini-0648:phyxminiproblems-problemphyxmini0648-positionintervalreadout}
  \lean{PhyXMiniProblems.ProblemPhyXMini0648.PositionIntervalReadout}
  A closed interval of numerical position readouts, measured in millimetres
\end{definition}
\begin{proof}
  This is the declared model definition of Position Interval Readout.
\end{proof}
\begin{definition}[region]
  \label{def:physics:phyx-mini-0648:phyxminiproblems-problemphyxmini0648-positionintervalreadout-region}
  \lean{PhyXMiniProblems.ProblemPhyXMini0648.PositionIntervalReadout.region}
  \uses{def:physics:phyx-mini-0648:phyxminiproblems-problemphyxmini0648-positionintervalreadout}
  The measurable position event represented by an interval readout
\end{definition}
\begin{proof}
  This is the declared model definition of region.
\end{proof}
\begin{definition}[interval Probability]
  \label{def:physics:phyx-mini-0648:phyxminiproblems-problemphyxmini0648-intervalprobability}
  \lean{PhyXMiniProblems.ProblemPhyXMini0648.intervalProbability}
  \uses{def:physics:phyx-mini-0648:phyxminiproblems-problemphyxmini0648-neutronpositionexperiment, def:physics:phyx-mini-0648:phyxminiproblems-problemphyxmini0648-positionintervalreadout}
  The dimensionless probability assigned to a position interval
\end{definition}
\begin{proof}
  This is the declared model definition of interval Probability.
\end{proof}
\begin{definition}[requested Central Interval]
  \label{def:physics:phyx-mini-0648:phyxminiproblems-problemphyxmini0648-requestedcentralinterval}
  \lean{PhyXMiniProblems.ProblemPhyXMini0648.requestedCentralInterval}
  \uses{def:physics:phyx-mini-0648:phyxminiproblems-problemphyxmini0648-positionintervalreadout}
  The interval from x = -1 mm through x = 1 mm asked for in the problem
\end{definition}
\begin{proof}
  This is the declared model definition of requested Central Interval.
\end{proof}
\begin{definition}[Answer Choice]
  \label{def:physics:phyx-mini-0648:phyxminiproblems-problemphyxmini0648-answerchoice}
  \lean{PhyXMiniProblems.ProblemPhyXMini0648.AnswerChoice}
  Labels of the four answer choices printed with the problem
\end{definition}
\begin{proof}
  This is the declared model definition of Answer Choice.
\end{proof}
\begin{definition}[probability Value]
  \label{def:physics:phyx-mini-0648:phyxminiproblems-problemphyxmini0648-answerchoice-probabilityvalue}
  \lean{PhyXMiniProblems.ProblemPhyXMini0648.AnswerChoice.probabilityValue}
  \uses{def:physics:phyx-mini-0648:phyxminiproblems-problemphyxmini0648-answerchoice}
  The dimensionless probability printed beside each answer choice
\end{definition}
\begin{proof}
  This is the declared model definition of probability Value.
\end{proof}
\begin{definition}[recorded Answer Choice]
  \label{def:physics:phyx-mini-0648:phyxminiproblems-problemphyxmini0648-recordedanswerchoice}
  \lean{PhyXMiniProblems.ProblemPhyXMini0648.recordedAnswerChoice}
  \uses{def:physics:phyx-mini-0648:phyxminiproblems-problemphyxmini0648-answerchoice}
  Dataset metadata records choice C; this constant is never a premise
\end{definition}
\begin{proof}
  This is the declared model definition of recorded Answer Choice.
\end{proof}
\begin{definition}[Matches Answer Choice]
  \label{def:physics:phyx-mini-0648:phyxminiproblems-problemphyxmini0648-matchesanswerchoice}
  \lean{PhyXMiniProblems.ProblemPhyXMini0648.MatchesAnswerChoice}
  \uses{def:physics:phyx-mini-0648:phyxminiproblems-problemphyxmini0648-neutronpositionexperiment, def:physics:phyx-mini-0648:phyxminiproblems-problemphyxmini0648-positionintervalreadout, def:physics:phyx-mini-0648:phyxminiproblems-problemphyxmini0648-intervalprobability, def:physics:phyx-mini-0648:phyxminiproblems-problemphyxmini0648-answerchoice}
  A choice matches when its displayed value equals the modeled probability
\end{definition}
\begin{proof}
  This is the declared model definition of Matches Answer Choice.
\end{proof}
\begin{theorem}[problem phyx mini 0648]
  \label{thm:physics:phyx-mini-0648:phyxminiproblems-problemphyxmini0648-problem-phyx-mini-0648}
  \lean{PhyXMiniProblems.ProblemPhyXMini0648.problem_phyx_mini_0648}
  \uses{def:physics:phyx-mini-0648:phyxminiproblems-problemphyxmini0648-neutronpositionexperiment, def:physics:phyx-mini-0648:phyxminiproblems-problemphyxmini0648-satisfiesonedimensionalpositionquantumlaws, def:physics:phyx-mini-0648:phyxminiproblems-problemphyxmini0648-matchesproblemstatementandfigure, def:physics:phyx-mini-0648:phyxminiproblems-problemphyxmini0648-intervalprobability, def:physics:phyx-mini-0648:phyxminiproblems-problemphyxmini0648-requestedcentralinterval}
  For the normalized two-plateau neutron state shown in the figure, the central 2 mm interval occupies one quarter of the constant-density 8 mm support
\end{theorem}
\begin{proof}
  The conclusion follows from the governing relations and figure readouts named in the statement of problem phyx mini 0648.
\end{proof}
% --- Archon declaration topology end ---
% --- Archon physics formalization source end ---
